\chapter{Workflows, Jobs, and Pipeline Orchestration}
\index{Workflows}\index{jobs}\index{tasks}\index{orchestration}\index{task values}%
\index{retry}\index{repair run}\index{cron}\index{if/else task}\index{for-each task}%
\begin{objectives}
\begin{itemize}
    \item Compose multi-task jobs with dependency DAGs, job clusters, and schedules.
    \item Pass data between tasks with job parameters and task values.
    \item Control flow with conditional (\texttt{if/else}) and \texttt{for each} task types.
    \item Engineer reliability: retries, timeouts, notifications, and repair runs.
    \item Orchestrate DLT pipelines, notebooks, and SQL in one governed job.
    \item Monitor an estate of jobs with system tables and define an on-call standard.
\end{itemize}
\end{objectives}

\begin{crosscloud}
Databricks Workflows is the platform's \textbf{native orchestrator}; the concepts---DAGs, tasks,
retries, SLAs---map onto every cloud's orchestrator.

\vspace{2mm}
{\footnotesize
\begin{tabular}{@{}p{3.0cm}p{3.4cm}p{3.2cm}p{3.2cm}@{}}
\toprule
\textbf{Capability} & \textbf{Databricks} & \textbf{AWS} & \textbf{Azure / GCP / Fabric} \\
\midrule
Orchestrator & Workflows (jobs + tasks) & Step Functions / MWAA (Airflow) & ADF / Fabric pipelines / Cloud Composer \\
Task compute & Job cluster per task/task group & Activity workers & Activity IRs / Composer workers \\
Parameters & Job params + task values & SFN input/output passing & Pipeline params / Airflow XComs \\
Conditional flow & if/else, for-each tasks & Choice / Map states & If Condition / ForEach activities \\
Partial recovery & Repair run (failed + downstream) & Redrive / Airflow retries & Rerun from activity / task retry \\
API/IaC & Jobs API, Asset Bundles & SFN ASL / CDK / Terraform & ARM / Bicep / Terraform \\
\bottomrule
\end{tabular}}

\vspace{1mm}
External orchestrators (Airflow, Dagster) still integrate via the Jobs API; the native advantage
is colocation with clusters, catalogs, and lineage---one pane for run, retry, and data impact.
\end{crosscloud}

\section{Week 12 Roadmap and Mental Model}
Chapters 9--11 built the pipeline \emph{components}; Week 12 assembles them into an \emph{operated system}. A Workflow job is a DAG of tasks---notebooks, DLT pipelines, SQL, JARs, dbt, conditionals---with scheduling, retries, notifications, and a repair model \cite{databricksWorkflows}. The engineering content is not drawing boxes; it is failure design: what happens at 03:00 when task 4 of 9 fails, and who finds out how.

The mental model: \textbf{a job is a small distributed system}. Tasks are processes with dependencies, identity (run-as principal), resources (job clusters), inputs (parameters, task values), and failure semantics (retry, timeout, downstream skip, repair). Design it like one and the 03:00 page is boring; design it like a script and the page is not.

\begin{definitionbox}
\textbf{A Databricks job} is a scheduled or triggered DAG of \textbf{tasks} executed by Workflows. Each task has a type (notebook, pipeline, SQL, condition, for-each), compute (job cluster, serverless, warehouse), dependencies (\texttt{depends\_on}), and reliability settings (retries, timeout). A \textbf{repair run} re-executes only failed tasks and their dependents.
\end{definitionbox}

\begin{keyterms}
\textbf{Job parameter}: a value set at job level, shared by all tasks. \textbf{Task value}: a runtime value set by one task (\texttt{dbutils.jobs.taskValues.set}) and read by dependents. \textbf{Job cluster}: ephemeral compute created for a task or run. \textbf{Repair run}: partial re-execution of failed and downstream tasks. \textbf{Run-as}: the identity the job executes as---a service principal in production.
\end{keyterms}

\section{Monday: Jobs, Tasks, and Compute Topology}
\subsection{Anatomy of a Job}
A job has: a schedule (cron or file-arrival/continuous trigger), tasks with \texttt{depends\_on} edges, per-task compute, notifications, and a run-as identity. The first design decision is compute topology: one shared job cluster for all tasks (cheap, coupled) versus a job cluster per task (isolated, slightly slower, failure-contained). Production standard in this book: \emph{per-task job clusters for anything that can fail independently}.

\subsection{Scheduling}
Cron schedules handle time-based runs; \textbf{file-arrival triggers} start jobs when data lands (composing beautifully with Chapter 11's partner files); continuous triggers keep jobs alive. Time zone and daylight-saving behavior are configuration, not folklore---set them explicitly.

\begin{pitfall}
\textbf{The shared-cluster cascade.} Six tasks on one job cluster: task 2's memory leak kills tasks 3--6 as ``failures'' that were really collateral damage. Per-task clusters cost a few extra start minutes and buy you accurate failure attribution---the repair run then retries the task that actually failed.
\end{pitfall}

\section{Tuesday: Parameters, Task Values, and Control Flow}
\subsection{Static versus Runtime Values}
\textbf{Job parameters} are declared at job level (e.g.\ \texttt{processing\_date}, \texttt{environment}) and referenced by every task. \textbf{Task values} are computed at runtime and passed downstream---the row count the ingestion task measured, the quality verdict the gate produced.

\begin{lstlisting}[language=Python,caption={Task values: the runtime handoff between tasks.},label={lst:ch12-taskvalues}]
# In the quality-gate task:
bad = spark.sql("""
    SELECT COUNT(*) c FROM de_training.silver.orders
    WHERE order_id IS NULL OR total_amount < 0
""").first().c

if bad > 0:
    dbutils.jobs.taskValues.set("quality", "failed")
    raise ValueError(f"Quality gate failed: {bad} bad rows")
dbutils.jobs.taskValues.set("quality", "passed")
dbutils.jobs.taskValues.set("rows_silver",
    spark.table("de_training.silver.orders").count())

# In a downstream task:
rows = dbutils.jobs.taskValues.get(task_key="quality_gate", key="rows_silver")
\end{lstlisting}

\subsection{Conditionals and For-Each}
An \texttt{if/else} task evaluates a condition (often against a task value) and routes the DAG---the in-band alerting and branching mechanism. A \texttt{for each} task fans out over a list (e.g.\ one sub-run per partner feed) with concurrency controls---the answer to ``we onboard 40 tables with one config table'' without 40 hand-written DAGs.

\section{Wednesday: Reliability Engineering}
\subsection{Retries, Timeouts, Notifications}
Retries handle transient failure (with backoff); timeouts bound hung tasks; notifications (email, webhook to PagerDuty/Teams/Slack) route failure to humans. The policy question is not ``retry or not'' but \emph{what is safe to retry}: idempotent tasks (Chapter 3's discipline) retry freely; side-effecting tasks retry only after a state check.

\subsection{Repair Runs}
When a run fails at task 4 of 9, a \textbf{repair run} re-executes task 4 and its dependents with the same parameters---the other five tasks' outputs are reused. Repair is the economic and operational alternative to full reruns, but it leans entirely on task idempotency: a non-idempotent task 2 ``reused'' by a repair is only safe because it succeeded the first time \cite{databricksJobsAPI}.

\subsection{Monitoring the Estate}
Job and run history live in \texttt{system.lakeflow} tables: success rates, duration trends, cost per job. The on-call standard this book adopts: every production job has an owner, a runbook link in its failure notification, and a weekly reliability review of the failure list.

\section{Thursday Hands-on Project: The Orchestrated Medallion}
Assemble Weeks 10--11 into one job: Auto Loader ingestion $\rightarrow$ DLT pipeline $\rightarrow$ quality-gate notebook $\rightarrow$ conditional alert, parameterized by date, with a forced failure and a repair run.
\index{hands-on lab}\index{repair run!lab}\index{Workflows!lab}

\begin{lstlisting}[language=json,caption={Week 12 lab: job definition (JSON view of the DAG).},label={lst:ch12-job}]
{
  "name": "de-week12-nightly-elt",
  "parameters": [{"name": "processing_date", "default": "{{ds}}"}],
  "tasks": [
    {"task_key": "ingest", "job_cluster_key": "jc",
     "notebook_task": {"notebook_path": "src/ingest_autoloader",
                       "base_parameters": {"processing_date": "{{job.parameters.processing_date}}"}}},
    {"task_key": "pipeline", "depends_on": [{"task_key": "ingest"}],
     "pipeline_task": {"pipeline_id": "<dlt-pipeline-id>"}},
    {"task_key": "quality_gate", "depends_on": [{"task_key": "pipeline"}],
     "job_cluster_key": "jc",
     "notebook_task": {"notebook_path": "src/quality_gate"}},
    {"task_key": "route", "depends_on": [{"task_key": "quality_gate"}],
     "condition_task": {"op": "EQUAL_TO",
                        "left": "{{tasks.quality_gate.values.quality}}",
                        "right": "passed"}},
    {"task_key": "alert", "depends_on": [{"task_key": "route", "outcome": "false"}],
     "notebook_task": {"notebook_path": "src/notify_oncall"}}
  ],
  "job_clusters": [{"job_cluster_key": "jc",
    "new_cluster": {"spark_version": "15.4.x-scala2.12",
                    "node_type_id": "Standard_DS3_v2", "num_workers": 1}}],
  "schedule": {"quartz_cron_expression": "0 0 2 * * ?", "timezone_id": "UTC"}
}
\end{lstlisting}

The drill: run the job; force the quality gate to fail (point it at a table with planted bad rows); observe the conditional alert fire; fix the data; then \texttt{databricks runs repair -{}-run-id <id>} and capture that only failed and dependent tasks re-executed.

\subsection*{Deliverables}
\begin{itemize}
    \item The job definition (UI export or JSON) with the five-task DAG.
    \item The forced-failure evidence: gate failure, alert task execution, notification.
    \item The repair-run page showing partial re-execution.
\end{itemize}

\subsection*{Acceptance Criteria}
\begin{itemize}
    \item A same-day rerun (after success) changes zero Gold row counts---end-to-end idempotence.
    \item The forced failure pages via the conditional path, with the runbook link present.
    \item The repair run re-executes only the failed task and its dependents.
    \item The job runs as a service principal with exactly the grants from Chapter 8's matrix.
\end{itemize}

\subsection*{Stretch Goals}
\begin{itemize}
    \item Add a \texttt{for each} task fanning ingestion over the nine Olist tables from a config list.
    \item Add a file-arrival trigger variant and compare its latency against the cron schedule.
    \item Emit per-task duration and row counts to a run-ledger Delta table and chart a week of runs.
\end{itemize}

\section{Friday Use Case: Sixty Pipelines, Three Teams, One 6 A.M. SLA}
A company runs 60 interdependent nightly pipelines across three teams with a shared 06:00 completion SLA. Today: cron-scheduled notebooks, failures discovered by customers, and monthly all-hands reruns. Design the orchestration layer.
\index{orchestration!design}\index{SLA}

\subsection{Requirements and Assumptions}
Cross-team dependencies exist (finance consumes logistics' Gold). Backfills happen weekly. Teams deploy independently. The platform team owns the on-call rotation and the SLA.

\begin{figure}[H]
    \centering
    \resizebox{\textwidth}{!}{%
    \begin{tikzpicture}[
        team/.style={rectangle, rounded corners, draw=primary, thick, fill=primary!4, minimum width=4.2cm, minimum height=2.6cm},
        svc/.style={rectangle, rounded corners, draw=primary, thick, fill=white, align=center, minimum width=1.7cm, minimum height=0.7cm, font=\tiny\bfseries},
        plat/.style={rectangle, rounded corners, draw=codegreen!60!black, thick, fill=green!7, align=center, minimum width=3.4cm, minimum height=0.9cm, font=\scriptsize\bfseries},
        flow/.style={-{Latex[length=2mm]}, draw=primary, thick}
    ]
        \node[team, label={[font=\scriptsize\bfseries\color{primary}]above:logistics jobs}] (log) at (0,0) {};
        \node[team, label={[font=\scriptsize\bfseries\color{primary}]above:finance jobs}] (fin) at (5.4,0) {};
        \node[team, label={[font=\scriptsize\bfseries\color{primary}]above:marketing jobs}] (mkt) at (10.8,0) {};
        \node[svc] at (-0.9,0.5) {ingest}; \node[svc] at (1.1,0.5) {dlt};
        \node[svc] at (-0.9,-0.6) {gate};  \node[svc] at (1.1,-0.6) {publish};
        \node[svc] at (4.5,0.5) {close};   \node[svc] at (6.5,0.5) {report};
        \node[svc] at (4.5,-0.6) {gate};   \node[svc] at (6.5,-0.6) {publish};
        \node[svc] at (9.9,0.5) {features}; \node[svc] at (11.9,0.5) {score};
        \node[svc] at (9.9,-0.6) {gate};   \node[svc] at (11.9,-0.6) {publish};
        \node[plat] (ledger) at (5.4,-2.6) {run ledger + SLA monitor\\(system.lakeflow + alerts)};
        \draw[flow] (log) -- node[above, font=\scriptsize]{contract tables} (fin);
        \draw[flow] (fin) -- (ledger);
        \draw[flow] (log) -- (ledger);
        \draw[flow] (mkt) -- (ledger);
    \end{tikzpicture}%
    }
    \caption{Federated orchestration: team-owned jobs, contract tables between teams, one platform-owned ledger and SLA monitor.}
    \label{fig:ch12-estate}
\end{figure}

\subsection{Architecture Decisions}
\textbf{Federated jobs, contractual interfaces.} Each team owns its Workflow jobs; cross-team dependencies are \emph{data contracts}, not task edges across team boundaries---finance's job starts from a file-arrival trigger or a contract-table freshness check, not from a dependency on logistics' internals. This is what lets teams deploy independently without 03:00 cross-team debugging.

\textbf{One ledger, one SLA monitor.} Every job writes run outcomes to a shared run-ledger table (task values make this cheap); the platform SLA monitor reads \texttt{system.lakeflow} plus the ledger, pages on trajectory-to-miss at 04:30 rather than on breach at 06:00, and drives the weekly reliability review. Backfills are bundle-parameterized runs (Chapter 23) against the same jobs, never hand-edited clones.

\begin{pitfall}
\textbf{Cross-team task dependencies.} Wiring team A's task directly into team B's job graph couples their deploys, retries, and failure pages. Depend on \emph{published data with a freshness contract}; alert the producer when the contract breaks.
\end{pitfall}

\subsection{Risks and Mitigations}
\begin{table}[H]
\centering
\small
\begin{tabular}{@{}p{4.6cm}p{7.6cm}@{}}
\toprule
\textbf{Risk} & \textbf{Mitigation} \\
\midrule
Cross-team dependency couples releases & Contract tables with freshness SLAs, not task edges \\
Repair run re-executes a side-effecting task & Idempotency review is a job-definition checklist item \\
Alert fatigue from flaky tasks & Failure budgets per job; flakiness is a reliability bug \\
Run ledger becomes unqueryable lore & Ledger schema versioned; weekly review uses saved queries \\
On-call pages without runbooks & Notification template requires the runbook link to save \\
\bottomrule
\end{tabular}
\caption{Week 12 scenario risks and mitigations.}
\label{tab:ch12-risks}
\end{table}

\section{Design Decision Guide}
\index{decision guide!Week 12}%
Orchestration decisions are failure-shape decisions: choose based on what 03:00 looks like.

\begin{table}[H]
\centering
\small
\begin{tabular}{@{}p{3.4cm}p{4.4cm}p{4.6cm}@{}}
\toprule
\textbf{Decision} & \textbf{Choose the first when} & \textbf{Choose the second when} \\
\midrule
Per-task vs.\ shared job cluster & Tasks can fail independently & Tiny chained tasks where start-up dominates \\
Task value vs.\ job parameter & Value is computed at runtime & Value is declared input \\
Conditional task vs.\ external alerting & The response is in-band and automatable & Humans must decide with context \\
Workflows vs.\ external orchestrator & The estate is Databricks-centric & Orchestration spans many external systems \\
Repair run vs.\ full rerun & Tasks are idempotent & State is ambiguous---then fix idempotency \\
Cron vs.\ file-arrival trigger & Time-driven SLAs & Data-arrival-driven work \\
\bottomrule
\end{tabular}
\caption{Week 12 decision guide.}
\label{tab:ch12-decisions}
\end{table}

Idempotency is the load-bearing assumption under repair runs, retries, and backfills. If you take one thing from Week 12: orchestration features are only as safe as the write patterns beneath them.

\section{Operational Guidance for Week 12}
\index{operational guidance!Week 12}%
\subsection{Performance and Reliability}
The run ledger is the estate's health record: review failure rate, mean repair time, and p95 duration per job weekly. Trajectory alerting (elapsed + remaining-p50 vs.\ SLA) pages early; breach alerts write postmortems. Keep repair drills on the calendar---a repair path exercised monthly works at 03:00; one last used at launch is a rumor.

\subsection{Security and Governance Checklist}
\begin{itemize}
    \item Production jobs run-as service principals with Chapter 8's least-privilege matrix.
    \item Notification webhooks (PagerDuty/Teams) carry no payload data---keys and links only.
    \item Job parameters containing sensitive values come from secret scopes, not job config.
    \item Schedule time zones set explicitly and reviewed before daylight-saving transitions.
\end{itemize}

\subsection{Troubleshooting Playbook}
\begin{table}[H]
\centering
\small
\begin{tabular}{@{}p{3.6cm}p{4.2cm}p{4.8cm}@{}}
\toprule
\textbf{Symptom} & \textbf{Likely cause} & \textbf{First action} \\
\midrule
Downstream task reads empty value & Task value key/task\_key mismatch & Log the set/get keys; check dependency edge \\
Condition task always false & Type mismatch in comparison & Compare string forms; log the evaluated value \\
Repair re-ran half the job & Over-declared dependencies & Tighten \texttt{depends\_on} to true data edges \\
Job fired an hour late/early & Time zone or DST in cron schedule & Set \texttt{timezone\_id} explicitly; re-check at DST \\
\bottomrule
\end{tabular}
\caption{Week 12 troubleshooting quick reference.}
\label{tab:ch12-trouble}
\end{table}

\section{Interview Q\&A}
\subsection{Concepts and Trade-offs}
\begin{enumerate}
    \item \textbf{Q: Task values versus job parameters?}\\
    A: Parameters are static inputs declared at job level; task values are runtime outputs of one task consumed by dependents. Confusing them is how ``yesterday's date'' ends up hardcoded.
    \item \textbf{Q: What makes a repair run safe?}\\
    A: Task idempotency plus correct dependency modeling: repair reuses successful outputs and re-runs the failed subtree, so each task must produce the same effect when executed twice with the same inputs.
    \item \textbf{Q: How do you alert on trajectory, not breach?}\\
    A: Track per-task duration baselines in a run ledger; at 04:30 compare elapsed + remaining-p50 against the SLA and page early. Breach alerts tell you what you already lost.
    \item \textbf{Q: Job cluster per task or one shared cluster?}\\
    A: Per task for failure isolation and accurate attribution; shared only for tight-latency chains of small tasks where start-up dominates. State the trade-off explicitly in the job's README.
    \item \textbf{Q: When do you reach for Airflow instead of Workflows?}\\
    A: When orchestration must span many non-Databricks systems with complex scheduling logic already expressed in Airflow---and even then, Workflows runs the Databricks tasks via the Jobs API while Airflow coordinates the perimeter.
    \item \textbf{Q: How do you onboard a new pipeline into an existing job estate?}\\
    A: As a new job with its own SLA and ledger entries, depending on contract tables---never as a task inside another team's job.
    \item \textbf{Q: What does the repair-run button assume about your code?}\\
    A: That every task is idempotent and dependencies reflect true data edges. Repair is safe exactly to the extent your write discipline is real.
    \item \textbf{Q: How would you decommission a job safely?}\\
    A: Lineage first (who reads its outputs), then a deprecation notice in the run ledger, then schedule removal, then data-retention handling for its tables---the reverse of onboarding, with the same discipline.
\end{enumerate}

\section{Further Reading}
The Workflows documentation covers jobs, task types, triggers, and repair \cite{databricksWorkflows}; the Jobs API is the automation surface for Chapter 23 \cite{databricksJobsAPI}. System-table monitoring is documented in \cite{databricksSystemTables}. Revisit DLT \cite{databricksDLT} for how pipelines appear as tasks.

\section*{Key Takeaways}
\begin{itemize}
    \item A job is a small distributed system: identity, compute, inputs, and failure semantics per task.
    \item Per-task job clusters buy failure isolation and honest attribution for a few start-up minutes.
    \item Task values carry runtime facts; parameters carry declared inputs---keep them separate.
    \item Repair runs are the economics of recovery; they rest entirely on idempotent tasks.
    \item Federate orchestration at team boundaries with data contracts and freshness checks, not cross-team task edges.
\end{itemize}

\section{Exercises}
\begin{enumerate}
    \item \textbf{(Conceptual)} Diagram the difference between retry, rerun, and repair; give a scenario where each is the right response.
    \item \textbf{(Conceptual)} Why does a shared job cluster corrupt failure attribution?
    \item \textbf{(Hands-on)} Complete the Thursday lab, including the forced failure and repair evidence.
    \item \textbf{(Hands-on)} Add the run-ledger emission to every task and produce the week-in-review query.
    \item \textbf{(Hands-on)} Convert the cron schedule to a file-arrival trigger; measure arrival-to-start latency.
    \item \textbf{(Design)} Write the data contract (schema, freshness SLA, quality gate) between logistics and finance from the Friday scenario.
    \item \textbf{(Operations)} Draft the 04:30 trajectory alert query against the run ledger.
    \item \textbf{(Architecture)} Which tasks of this job become bundle resources in Chapter 23, and which stay UI-managed? Defend the split.
\end{enumerate}

\section{Milestone 3 Capstone: The Orchestrated Batch ELT Platform}\label{sec:ch12-capstone}
\index{Milestone 3 capstone}\index{capstone project!Part III}\index{Workflows!capstone}\index{Auto Loader!capstone}

\begin{capstonebox}
\textbf{Mission.} Close Part III by delivering the batch platform as one operated system: Auto Loader ingestion of daily Olist extracts, a DLT pipeline with expectations into Silver/Gold, a blocking quality gate, conditional failure alerting, and a demonstrated repair run---all parameterized and idempotent.

\textbf{Estimated time.} 6--8 hours.

\textbf{What you'll practice.}
\begin{itemize}
    \item Triggered Auto Loader with schema evolution and rescued data.
    \item DLT expectations wired to real severity decisions.
    \item Multi-task Workflow design with task values, conditionals, and repair.
    \item Idempotence proof and DBU cost accounting per nightly run.
\end{itemize}
\end{capstonebox}

\subsection*{Scenario}
Contoso Retail's nightly load has grown from one CSV to nine partner feeds plus a CDC extract. The current cron-and-notebook setup double-loads on reruns and fails silently twice a month. You are replacing it with the orchestrated platform---and proving every failure mode is now loud, repairable, and reconciled.

\subsection*{Requirements}
\begin{itemize}
    \item One Workflow job: ingest (Auto Loader) $\rightarrow$ DLT medallion pipeline $\rightarrow$ quality gate $\rightarrow$ publish $\rightarrow$ conditional alert path.
    \item Schema-evolving batch and one corrupted row handled without manual repair.
    \item Expectations with deliberate severities; quality metrics queryable from the event log.
    \item Forced mid-pipeline failure repaired (not fully rerun) with evidence.
    \item Same-day idempotence proof; per-run DBU cost from system tables.
    \item README with architecture diagram, runbook, and failure drill results.
\end{itemize}

\subsection*{Reference Architecture}
Reuse \cref{fig:ch12-estate}'s single-team slice: ingestion task $\rightarrow$ DLT task $\rightarrow$ gate $\rightarrow$ publish, with the conditional alert branch and the run ledger---deployed as one job with per-task job clusters.

\subsection*{Implementation Steps}
\begin{enumerate}
    \item Stand up the Auto Loader volume ingestion with checkpoint and schema location under governance.
    \item Convert the medallion flow to DLT with expectations (reuse Chapter 10).
    \item Author the job JSON per \cref{lst:ch12-job}; wire task values and the conditional alert.
    \item Run the failure drill; repair; capture evidence.
    \item Produce the cost report and the runbook; finalize the README.
\end{enumerate}

\subsection*{Deliverables}
\begin{itemize}
    \item Job definition, DLT pipeline source, and the quality-gate notebook in Git.
    \item Failure-drill and repair-run evidence; idempotence proof.
    \item Cost-per-run report with the job-cluster decision documented.
\end{itemize}

\subsection*{Acceptance Criteria}
\begin{itemize}
    \item A poisoned feed demonstrably drops/fails per expectation severity, with event-log metrics.
    \item The repair run re-executes only the failed subtree; total rerun cost is documented.
    \item Gold row counts are identical across a forced same-day rerun.
    \item The job runs as a service principal; no credentials in any committed file.
\end{itemize}

\subsection*{Stretch Goals}
\begin{itemize}
    \item Add the \texttt{for each} fan-out over the nine feeds from a config table.
    \item Add the trajectory-based SLA alert from the Friday scenario.
    \item Wrap the job in an Asset Bundle with dev/prod targets (preview of Chapter 23).
\end{itemize}
